\subsection{折叠群胚矩阵单位对称、正常化子交换化与分层奇偶证书}\label{subsec:conclusion-fold-groupoid-matrixunit-normalizer-parity}
本小节把折叠等价关系群胚的离散 \(\ast\)-对称、纤维正常化子 \(N_m\) 的花环骨架以及由纤维规模分层产生的全部 \(\{\pm1\}\)-证书统一到同一结论链中。
附录中的群胚代数模型给出“矩阵单位基置换型 \(\ast\)-自同构”与可观测提升群的严格等同，而花环分解则把交换化、交换子与指数 \(2\) 商完全压缩为有限个可审计奇偶字符。

\begin{theorem}[群胚矩阵单位对称的正常化子刻画]\label{thm:conclusion-fold-groupoid-matrixunit-normalizer}
\leanverified{paper\_conclusion\_fold\_groupoid\_matrixunit\_normalizer}
沿用附录定理 \ref{thm:op-algebra-fold-quantum-channel-weak-hopf-groupoid} 的等价关系群胚
\[
\mathcal G_m=\Omega_m\times_{X_m}\Omega_m
\]
及其群胚代数 \(\CC[\mathcal G_m]\)，并定义
\[
\Aut_{\mathrm{MU}}^{\ast}\!\bigl(\CC[\mathcal G_m]\bigr)
:=
\left\{
\Phi\in\Aut^{\ast}\!\bigl(\CC[\mathcal G_m]\bigr):
\Phi\ \text{将矩阵单位基}\ \{e_{\omega,\omega'}\}_{(\omega,\omega')\in\mathcal G_m}\ \text{作置换}
\right\}.
\]
则存在自然群同构
\[
N_m\ \cong\ \Aut_{\mathrm{MU}}^{\ast}\!\bigl(\CC[\mathcal G_m]\bigr),
\qquad
\tau\longmapsto \Phi_\tau,
\]
其中
\[
\Phi_\tau(e_{\omega,\omega'})=e_{\tau(\omega),\tau(\omega')}.
\]
\end{theorem}
\begin{proof}
这正是附录定理 \ref{thm:op-algebra-fold-groupoid-matrix-unit-aut-normalizer} 的结论。证毕。
\end{proof}

\begin{theorem}[正常化子交换化与全部 \texorpdfstring{$\{\pm1\}$}{+-1} 特征的分层完备分类]\label{thm:conclusion-fold-normalizer-abelianization-layered-characters}
沿用附录定理 \ref{thm:op-algebra-fold-fiber-normalizer-wreath} 的记号，并记
\[
X_m^{(d)}:=\{x\in X_m:\ d_m(x)=d\},
\qquad
N_{m,d}:=\card{X_m^{(d)}}.
\]
令
\[
\alpha_m
:=
\sum_{d\ge 2}\mathbf 1_{\{N_{m,d}\ge 1\}}
\;+\;
\sum_{d\ge 1}\mathbf 1_{\{N_{m,d}\ge 2\}}.
\]
对 \(\tau\in N_m\)，记 \(f:=\pi(\tau)\in \Sym_{d_m}(X_m)\) 为其在 \(X_m\) 上诱导的置换。
任取附录命题 \ref{prop:op-algebra-fold-lift-global-sign-factorization} 中的纤维标号，并记
\[
p_x(\tau)\in \mathfrak S_{d_m(x)}
\]
为 \(\tau\) 在第 \(x\) 条纤维上诱导的局部置换。
则对每个满足 \(d\ge 2\) 且 \(N_{m,d}\ge 1\) 的 \(d\)，量
\[
\varepsilon_d(\tau)
:=
\prod_{x\in X_m^{(d)}}\operatorname{sgn}\bigl(p_x(\tau)\bigr)\in\{\pm1\}
\]
定义一个群同态；对每个满足 \(N_{m,d}\ge 2\) 的 \(d\)，量
\[
\delta_d(\tau)
:=
\operatorname{sgn}\bigl(f|_{X_m^{(d)}}\bigr)\in\{\pm1\}
\]
亦定义一个群同态。
并且有自然同构
\[
N_m^{\mathrm{ab}}
:=
N_m/[N_m,N_m]
\ \cong\
(\ZZ/2\ZZ)^{\alpha_m},
\]
以及
\[
\Hom(N_m,\{\pm1\})
\ \cong\
(\ZZ/2\ZZ)^{\alpha_m},
\]
其中一组标准基恰可取为
\[
\{\varepsilon_d:\ d\ge 2,\ N_{m,d}\ge 1\}
\ \cup\
\{\delta_d:\ d\ge 1,\ N_{m,d}\ge 2\}.
\]
\end{theorem}
\begin{proof}
附录定理 \ref{thm:op-algebra-fold-wreath-abelianization-characters} 已给出 \(N_m\) 的全部 \(\{\pm1\}\)-角色分类。
其中
\[
\varepsilon_d=\chi_d^{\mathrm{fib}},
\qquad
\delta_d=\chi_d^{\mathrm{vis}},
\]
故所列角色正是该定理中的显式基。
再按存在性条件逐个计数其基向量个数，即得到
\(
\alpha_m
\)
与上述两个 \((\ZZ/2\ZZ)^{\alpha_m}\) 同构。证毕。
\end{proof}

\begin{theorem}[折叠正常化子的一维特征完全由纤维规模分层决定]\label{thm:conclusion-fold-normalizer-characters-completely-stratified}
沿用定理 \ref{thm:conclusion-fold-normalizer-abelianization-layered-characters} 的记号。
则存在规范满同态
\[
\pi_m:N_m\twoheadrightarrow (\ZZ/2\ZZ)^{\alpha_m}
\]
使得任意 \(\{\pm1\}\)-值特征
\[
\chi:N_m\to\{\pm1\}
\]
都唯一经由 \(\pi_m\) 因子化。
等价地，
\[
\Hom(N_m,\{\pm1\})\cong (\ZZ/2\ZZ)^{\alpha_m},
\]
并且这一一维角色代数完全由纤维规模分层
\(
\{N_{m,d}\}_{d\ge 1}
\)
决定，而与各纤维内部更细几何无关。
\end{theorem}
\begin{proof}
定理 \ref{thm:conclusion-fold-normalizer-abelianization-layered-characters} 已给出
\[
N_m^{\mathrm{ab}}\cong (\ZZ/2\ZZ)^{\alpha_m},
\qquad
\Hom(N_m,\{\pm1\})\cong (\ZZ/2\ZZ)^{\alpha_m},
\]
其中一组标准基正由分层字符
\(
\{\varepsilon_d\}\cup\{\delta_d\}
\)
给出。
取交换化商与上述同构的复合即可定义 \(\pi_m\)。
由于任意 \(\{\pm1\}\)-值特征都必须经由交换化因子化，且该交换化已被这些分层奇偶字符完全张成，故因子化唯一。
最后，\(\alpha_m\) 与对应的基字符是否出现只取决于各层计数 \(N_{m,d}\)，从而全部一维特征完全由纤维规模分层决定。证毕。
\end{proof}

\begin{corollary}[分层奇偶可见量的统计盲性]\label{cor:conclusion-fold-normalizer-layered-character-histogram-blindness}
设 \(f:\Omega_f\twoheadrightarrow X_f\) 与 \(g:\Omega_g\twoheadrightarrow X_g\) 为两组有限折叠，其对应正常化子分别记为 \(N_f,N_g\)，并记
\[
X_f^{(d)}:=\{x\in X_f:\ \card{f^{-1}(x)}=d\},
\qquad
X_g^{(d)}:=\{y\in X_g:\ \card{g^{-1}(y)}=d\}.
\]
若对一切 \(d\ge 1\) 都有
\[
\card{X_f^{(d)}}=\card{X_g^{(d)}},
\]
则
\[
\Hom(N_f,\{\pm1\})\cong \Hom(N_g,\{\pm1\})
\]
作为由纤维规模分层标记的特征系完全一致。
因此，正常化子的一维 parity 审计只能看见规模层直方图，而看不见同规模层内部的任何非平凡排列几何。
\end{corollary}
\begin{proof}
把定理 \ref{thm:conclusion-fold-normalizer-characters-completely-stratified} 分别应用于 \(f\) 与 \(g\)。
它们的一维特征代数都只由各尺寸层是否出现以及是否出现至少两次决定，亦即只由分层计数
\(
\{\card{X_f^{(d)}}\}_{d\ge 1}
\)
与
\(
\{\card{X_g^{(d)}}\}_{d\ge 1}
\)
决定。
若这两组数据逐层相同，则对应的 \(\alpha\)-参数、分层奇偶基以及全部 \(\{\pm1\}\)-值特征系统都完全一致，故结论成立。证毕。
\end{proof}

\begin{proposition}[交换子群作为全部分层奇偶证书的公共核]\label{prop:conclusion-fold-normalizer-commutator-kernel}
沿用定理 \ref{thm:conclusion-fold-normalizer-abelianization-layered-characters} 的记号，则
\[
[N_m,N_m]
=
\bigcap_{\substack{d\ge 2\\ N_{m,d}\ge 1}}\ker(\varepsilon_d)
\ \cap\
\bigcap_{\substack{d\ge 1\\ N_{m,d}\ge 2}}\ker(\delta_d).
\]
更具体地，对每个 \(d\ge 2\)、\(n\ge 2\)，若记
\[
W_{d,n}:=\mathfrak S_d\wr \mathfrak S_n=\mathfrak S_d^{\,n}\rtimes \mathfrak S_n,
\]
则其交换子群满足
\[
[W_{d,n},W_{d,n}]
=
\Bigl\{
(g_1,\dots,g_n;\sigma)\in \mathfrak S_d^{\,n}\rtimes \mathfrak S_n:
\ \sigma\in A_n,\ 
\prod_{i=1}^{n}\operatorname{sgn}(g_i)=+1
\Bigr\}.
\]
因此 \([N_m,N_m]\) 由各 \(d\)-层上的上述局部交换子核作直积给出。
\end{proposition}
\begin{proof}
由定理 \ref{thm:conclusion-fold-normalizer-abelianization-layered-characters}，映射
\[
\Theta_m:N_m\to (\ZZ/2\ZZ)^{\alpha_m},
\qquad
\Theta_m(\tau):=\bigl(\{\varepsilon_d(\tau)\},\{\delta_d(\tau)\}\bigr)
\]
给出 \(N_m\) 的交换化满射。
故
\[
[N_m,N_m]=\ker(\Theta_m),
\]
即为全部分层奇偶证书的公共核。

对局部花环因子 \(W_{d,n}\)（\(d\ge 2,\ n\ge 2\)），附录定理 \ref{thm:op-algebra-fold-wreath-abelianization-characters} 表明其交换化恰由两条 \(\ZZ/2\ZZ\)-字符生成：一条记录基群 \(\mathfrak S_d^{\,n}\) 的总符号，另一条记录可见层 \(\mathfrak S_n\) 的符号。
因此交换子群就是这两条字符的公共核，亦即所示集合。
最后结合附录定理 \ref{thm:op-algebra-fold-fiber-normalizer-wreath} 的花环直积分解，即得 \(N_m\) 的结论。证毕。
\end{proof}

\begin{corollary}[全部指数 \texorpdfstring{$2$}{2} 商与指数 \texorpdfstring{$2$}{2} 正规子群的计数闭式]\label{cor:conclusion-fold-normalizer-index-two-quotients}
沿用定理 \ref{thm:conclusion-fold-normalizer-abelianization-layered-characters} 的记号，则 \(N_m\) 的指数 \(2\) 正规子群总数满足
\[
\#\{H\lhd N_m:\ [N_m:H]=2\}=2^{\alpha_m}-1.
\]
并且每个指数 \(2\) 正规子群都唯一写成某个非零角色
\[
\prod_{\substack{d\ge 2\\ N_{m,d}\ge 1}}\varepsilon_d^{\,a_d}\cdot
\prod_{\substack{d\ge 1\\ N_{m,d}\ge 2}}\delta_d^{\,b_d},
\qquad
a_d,b_d\in\{0,1\},
\]
的核。
\end{corollary}
\begin{proof}
指数 \(2\) 子群必为正规子群，并与满射
\(
N_m\twoheadrightarrow \ZZ/2\ZZ
\)
的核一一对应。
由定理 \ref{thm:conclusion-fold-normalizer-abelianization-layered-characters}，
\(
\Hom(N_m,\{\pm1\})\cong (\ZZ/2\ZZ)^{\alpha_m}
\)，
故非零角色恰有 \(2^{\alpha_m}-1\) 个。
标准基给出其唯一的 \(\ZZ_2\)-线性展开，于是得到所述唯一表示。证毕。
\end{proof}

\begin{theorem}[全局符号同态的纤维规模分层精确因子化]\label{thm:conclusion-fold-normalizer-global-sign-layer-factorization}
\leanverified{paper\_conclusion\_fold\_normalizer\_global\_sign\_layer\_factorization}
沿用定理 \ref{thm:conclusion-fold-normalizer-abelianization-layered-characters} 的记号。
则对任意 \(\tau\in N_m\) 有
\[
\operatorname{sgn}(\tau)
=
\Bigl(\prod_{\substack{d\ \mathrm{odd}\\ N_{m,d}\ge 2}}\delta_d(\tau)\Bigr)\cdot
\Bigl(\prod_{\substack{d\ge 2\\ N_{m,d}\ge 1}}\varepsilon_d(\tau)\Bigr).
\]
等价地，全局符号在上述角色基下由“奇纤维层的可见搬运奇偶”与“全部非平凡纤维层的内部奇偶”唯一决定。
\end{theorem}
\begin{proof}
这正是附录推论 \ref{cor:op-algebra-fold-global-sign-unique-decomposition} 在记号
\(
\delta_d=\chi_d^{\mathrm{vis}}
\),
\(
\varepsilon_d=\chi_d^{\mathrm{fib}}
\)
下的重写。
当 \(N_{m,d}\le 1\) 时，\(\delta_d\) 为平凡字符，故仅保留上式中出现的因子即可。证毕。
\end{proof}

\begin{conjecture}[正常化子群的分辨率刚性]\label{conj:conclusion-fold-normalizer-resolution-rigidity}
存在 \(m_0\ge 1\)，使得对任意 \(m,m'\ge m_0\)，若
\[
N_m\cong N_{m'}
\]
作为抽象群同构，则必有
\[
\bigl(N_{m,d}\bigr)_{d\ge 1}=\bigl(N_{m',d}\bigr)_{d\ge 1},
\]
从而 \(m=m'\)。
等价地，在大分辨率范围内，折叠群胚矩阵单位基的置换型 \(\ast\)-对称骨架
\[
\Aut_{\mathrm{MU}}^{\ast}\!\bigl(\CC[\mathcal G_m]\bigr)
\]
应当唯一决定分辨率及其纤维直方图。
\end{conjecture}
